\documentclass[tikz,border=6pt]{standalone}
\usepackage{pgfplots}
\pgfplotsset{compat=1.18}
\usepackage{amsmath}

\begin{document}
\begin{tikzpicture}
\begin{axis}[
    width=11cm, height=6.5cm,
    axis lines=left,
    axis line style={-latex},
    xlabel={}, ylabel={},
    xmin=0, xmax=8.6,
    ymin=0, ymax=1.08,
    xtick={2.0,4.2,5.1,6.2},
    xticklabels={
        {$\vphantom{q_1'}q_1$},
        {$q_1'$},
        {$\vphantom{q_1'}q_2$},
        {$\vphantom{q_1'}q_3$}
    },
    ytick={0.32,0.52,0.78},
    yticklabels={$p_1$,$p_2$,$p_3$},
    tick align=outside,
    clip=false,
]

% Quantile lookup guides.
\draw[densely dashed, green!45!black]
    (axis cs:0,0.32) -- (axis cs:2.0,0.32)
    (axis cs:2.0,0) -- (axis cs:2.0,0.32);
\draw[densely dashed, green!45!black]
    (axis cs:0,0.52) -- (axis cs:5.1,0.52)
    (axis cs:5.1,0) -- (axis cs:5.1,0.52);
\draw[densely dashed, green!45!black]
    (axis cs:0,0.78) -- (axis cs:6.2,0.78)
    (axis cs:6.2,0) -- (axis cs:6.2,0.78);

% A continuous rise to p_1, followed by a flat stretch.
\addplot[thick, blue!35!black, smooth]
    coordinates {(0,0) (0.35,0.08) (0.85,0.18) (1.4,0.27) (2.0,0.32)};
\addplot[thick, red!75!black]
    coordinates {(2.0,0.32) (4.2,0.32)};

% The CDF rises again and then jumps at q_3.
\addplot[thick, blue!35!black, smooth]
    coordinates {(4.2,0.32) (4.55,0.38) (5.1,0.52) (5.65,0.63) (6.2,0.70)};
\addplot[thick, blue!35!black, smooth]
    coordinates {(6.2,0.88) (6.75,0.94) (7.45,0.98) (8.6,1.01)};

% Right-continuity at the jump: the lower endpoint is excluded and the
% upper endpoint is included.
\node[circle, draw=blue!35!black, fill=white, thick, inner sep=2pt]
    at (axis cs:6.2,0.70) {};
\node[circle, fill=blue!35!black, inner sep=2.4pt]
    at (axis cs:6.2,0.88) {};

% Highlight the endpoints of the flat stretch and the attained p_2 lookup.
\node[circle, fill=blue!35!black, inner sep=1.8pt]
    at (axis cs:2.0,0.32) {};
\node[circle, fill=blue!35!black, inner sep=1.8pt]
    at (axis cs:4.2,0.32) {};
\node[circle, fill=blue!35!black, inner sep=1.8pt]
    at (axis cs:5.1,0.52) {};

\node[blue!35!black, anchor=west] at (axis cs:7.25,0.88) {$F_X(x)$};

\end{axis}
\end{tikzpicture}
\end{document}
